\documentclass[11pt,a4paper]{article}
\usepackage[T1]{fontenc}
\usepackage{lmodern,microtype,amsmath,amssymb,amsthm,mathtools}
\usepackage[margin=25mm,headheight=14pt]{geometry}
\usepackage{booktabs,array,tabularx,enumitem,xcolor,fancyhdr,listings,hyperref}
\definecolor{ink}{HTML}{17324D}
\definecolor{accent}{HTML}{176B78}
\definecolor{codebg}{HTML}{F5F7F9}
\hypersetup{colorlinks=true,linkcolor=ink,citecolor=accent,urlcolor=accent,
pdftitle={A constructive proof of the quintic alpha+n conjecture},
pdfsubject={Positive integral palette lifting and exact quintic chromatic factors}}
\newtheorem{theorem}{Theorem}[section]
\newtheorem{lemma}[theorem]{Lemma}
\newtheorem{proposition}[theorem]{Proposition}
\newtheorem{corollary}[theorem]{Corollary}
\theoremstyle{definition}
\newtheorem{remark}[theorem]{Remark}
\newcommand{\Z}{\mathbb Z}
\newcommand{\Q}{\mathbb Q}
\newcommand{\R}{\mathbb R}
\newcommand{\fall}[2]{(#1)_{#2}}
\newcommand{\htf}{\mathcal H}
\DeclareMathOperator{\diag}{diag}
\pagestyle{fancy}
\fancyhf{}
\fancyhead[L]{\small\color{ink}Quintic chromatic factors}
\fancyhead[R]{\small 28 November 2026}
\fancyfoot[C]{\small\thepage}
\setlength{\emergencystretch}{3em}
\setlist{nosep,leftmargin=*}
\raggedbottom
\lstset{language=Python,basicstyle=\ttfamily\scriptsize,
keywordstyle=\color{ink}\bfseries,commentstyle=\color{accent},
backgroundcolor=\color{codebg},frame=single,rulecolor=\color{black!12},
breaklines=true,columns=fullflexible,keepspaces=true,
showstringspaces=false,tabsize=4,aboveskip=10pt,belowskip=10pt}
\title{\color{ink}\bfseries A constructive proof of the quintic\\
\texorpdfstring{\boldmath $\alpha+n$}{alpha+n} conjecture}
\author{\normalsize Research manuscript}
\date{28 November 2026}

\begin{document}
\maketitle
\begin{abstract}
For every monic quintic $f\in\Z[x]$, we construct positive integers $N,n$
and a connected finite simple graph $G$ such that
$P_G(q)=\fall qN f(q-n)$.
The graph consists of an $N$-clique and a five-clique with specified
edges between them. The construction starts from thirty-one strictly
positive rational palette atoms whose leading matching polynomial is $X^5$.
Changing only intersections containing one distinguished row makes the
coefficient map jointly affine. An explicit matrix of determinant one
modulo $24$, followed by corrections in $12\Z$, $4\Z$, and $\Z$,
then gives an exact integral solution. Uniform bounds prove that all atoms
remain positive at the predetermined scale
$t=96000\htf$, where
$\htf=10+|a|+|b|+|c|+|d|+|e|$ for
$f=x^5+ax^4+bx^3+cx^2+dx+e$.
In particular, $|V(G)|\le18\,432\,000\,000\,\htf^2$.
The argument is self-contained and applies to every quintic, including
reducible polynomials. It implies the $\alpha+n$ statement for every
algebraic integer of degree at most five and an exact splitting-field
realization result in these degrees. A standard-library implementation
and an independent matching verification are included.
\end{abstract}

\section{The named problem and the result}

A chromatic root is a complex zero of the chromatic polynomial $P_G$ of
a finite simple graph $G$. Cameron and Morgan's $\alpha+n$ conjecture asks
whether, for every algebraic integer $\alpha$, some natural number $n$
makes $\alpha+n$ a chromatic root \cite[Conjecture 4]{CM}.
The quadratic case was established by the original study group, and Bohn
established the cubic case \cite{CM,Bohn}. We address the unrestricted
degree-five case of that same statement.

For $N\ge0$, write $\fall qN=q(q-1)\cdots(q-N+1)$, with $\fall q0=1$.

\begin{theorem}[Exact quintic realization]\label{thm:main}
Let
\[
 f(x)=x^5+ax^4+bx^3+cx^2+dx+e,\qquad a,b,c,d,e\in\Z.
\]
There are positive integers $N,n$ and a connected finite simple graph $G$
with $N+5$ vertices such that
\begin{equation}\label{eq:main}
 P_G(q)=\fall qN f(q-n).
\end{equation}
The graph is the union of cliques of orders $N$ and $5$, with specified
cross-edges, and its chromatic number is $N$.
Set
\begin{equation}\label{eq:scale}
 \htf=10+|a|+|b|+|c|+|d|+|e|,\qquad t=96000\htf.
\end{equation}
The explicit construction below works at this scale and satisfies
\begin{equation}\label{eq:bounds}
 0<n<t^2,\qquad N+5\le2t^2
 =18\,432\,000\,000\,\htf^2.
\end{equation}
No irreducibility, coefficient-sign, signature, or Galois-group assumption
is imposed.
\end{theorem}

\begin{corollary}\label{cor:degree}
Every algebraic integer of degree at most five has a positive integer
translate that is a chromatic root.
\end{corollary}
\begin{proof}
For a monic minimal polynomial $g$ of degree $j\le5$, apply
Theorem~\ref{thm:main} to $x^{5-j}g(x)$. If $g(\alpha)=0$, then
$P_G(\alpha+n)=0$.
\end{proof}

\subsection{Prior work and the role of the auxiliary problem}

The complement-of-a-bipartite-graph class used here is called a
$(5,N)$-\emph{biclique} in the cited literature; it is not a complete
bipartite graph. Partition and matching formulas for bicliques are
established tools \cite{Bohn,CM}. Gilmore's thesis gives computational
evidence for quartic and quintic fields and for possible quintic
signature/Galois-group types \cite[Chapter 7]{Gilmore}. Those statements
do not assert that every prescribed quintic algebraic integer has the
required integer translate. Delbourgo and Morgan's 2019 abstract describes
an algorithm for degrees at most four, conditional on existence of an
appropriate biclique \cite{DM19}. Their 2022 paper computes a fifth
chromatic coefficient and discusses the earlier evidence
\cite[Sections 1.1 and 1.3]{DM22}.

The auxiliary problem solved here is to prescribe all five coefficients
while retaining thirty-one positive integer palette atoms. The partition
identity alone does not solve this realization problem. The substantive
construction consists of the rational interior point in
Section~\ref{sec:profile}, the congruence matrix and integral corrections in
Sections~\ref{sec:matrices}--\ref{sec:lift}, and the uniform positivity
bounds in Section~\ref{sec:effective}. The affine observation is proved
directly and is not asserted, by itself, to be historically new.

Earlier quartic work in this research programme motivated the choice of a
positive quadratic margin. The supplied VMRC manuscript 245 likewise uses
descending coefficient correction with controlled lower-coordinate costs
\cite[Sections 2--3]{VMRC}. No theorem from either working manuscript is
assumed here. The present proof, including the case of degree at most four,
is self-contained. The scope of the literature review and its priority
limitation are stated in Section~\ref{sec:scope}.

\section{Palette atoms and affine dependence}\label{sec:palettes}

Let $C$ be an $N$-clique, and let $v_1,\ldots,v_5$ form a disjoint
five-clique. For each $i$, let $A_i\subseteq C$ be the vertices not
adjacent to $v_i$; every other cross-pair is an edge.
A vertex of type $\varnothing\ne J\subseteq[5]$ belongs to precisely
those $A_i$ with $i\in J$. Denote the number of such vertices by $m_J$.
We use no vertices of empty type.

For each nonempty $B\subseteq[5]$, prescribe an intersection moment $H_B$.
The corresponding atom counts are
\begin{equation}\label{eq:mobius}
 m_J=\sum_{B\supseteq J}(-1)^{|B|-|J|}H_B.
\end{equation}
If these thirty-one integers are nonnegative, they define actual sets:
create $m_J$ distinct vertices of each type. In particular, there is no
further set-realizability condition. Their union has size
\begin{equation}\label{eq:N}
 N=\sum_{\varnothing\ne B\subseteq[5]}(-1)^{|B|+1}H_B
   =\sum_{\varnothing\ne J\subseteq[5]}m_J.
\end{equation}

Fix an integer $L$, and write the moments as
\begin{equation}\label{eq:momentparams}
 H_{\{i\}}=L+z_i,\qquad H_B=L-D_B\quad (|B|\ge2).
\end{equation}
For a block $B$, let $c_{\{i\}}=z_i$, let $c_B=-D_B$ for $|B|\ge2$,
and let $\mu(B)=(-1)^{|B|-1}(|B|-1)!$.
For $R\subseteq[5]$, define
\begin{equation}\label{eq:FR}
 F_R(x)=\sum_{\pi\text{ a partition of }R}
          \prod_{B\in\pi}\mu(B)(x+c_B),\qquad F_\varnothing(x)=1.
\end{equation}
Write $F=F_{[5]}$.

\begin{lemma}[Exact counting]\label{lem:count}
When the moments are realized by nonnegative atoms,
\begin{equation}\label{eq:count}
 P_G(q)=\fall qN F(q-N+L).
\end{equation}
Moreover,
\begin{equation}\label{eq:trace}
 [x^4]F=\sum_{i=1}^5z_i-10.
\end{equation}
\end{lemma}
\begin{proof}
Color $C$ first with distinct colors and put $y=q-N$.
A collection of small-clique vertices indexed by $B$ has $y+H_B$
common available colors. With $x=y+L$, this is $x+c_B$.
Inclusion--exclusion on the partition lattice counts assignments in which
the five colors are all distinct, giving \eqref{eq:FR}. This is the standard
biclique partition formula specialized to five rows
\cite[Proposition 9]{CM}. For every sufficiently large integer $q$ the
count is valid, so it is a polynomial identity. The $N$-clique contributes
$\fall qN$.
The degree-four term comes from the singleton partition and the ten
partitions with one pair and three singletons, proving \eqref{eq:trace}.
\end{proof}

\begin{lemma}[One distinguished row]\label{lem:affine}
Hold all parameters on subsets of $[4]$ fixed, and distinguish row $5$.
The polynomial $F$ is jointly affine in $z_5$ and in all $D_B$ with
$5\in B$. For $5\in B$ and $|B|\ge2$,
\begin{equation}\label{eq:derivative}
 \frac{\partial F}{\partial D_B}
 =(-1)^{|B|}(|B|-1)!\,F_{[5]\setminus B}(x).
\end{equation}
Every polynomial on the right uses only the fixed outer rows.
\end{lemma}
\begin{proof}
Every set partition has exactly one block containing $5$.
Consequently no product in \eqref{eq:FR} contains two of the variable
parameters. Differentiating the factor of the block $B$ leaves its
Möbius coefficient with a minus sign and the sum over all partitions of
the complement, which is \eqref{eq:derivative}.
\end{proof}

For blocks of sizes $2,3,4,5$, the multipliers in
\eqref{eq:derivative} are $1,-2,6,-24$, and the complementary palette
polynomials have degrees $3,2,1,0$. The effects of successive corrections
add exactly by Lemma~\ref{lem:affine}.

\section{A strictly positive rational profile}\label{sec:profile}

For every $J\subseteq[5]$, write $\sigma_i(J)=1$ if $i\in J$ and
$\sigma_i(J)=-1$ otherwise. For nonempty $J$, set
\begin{equation}\label{eq:weights}
 w_J=\frac1{16}
 +\frac{\sigma_5(J)\bigl(\sigma_1(J)+\sigma_4(J)
                         -\sigma_2(J)-\sigma_3(J)\bigr)}{192}.
\end{equation}
The expression in parentheses has absolute value at most $4$, so
\begin{equation}\label{eq:wpositive}
 \frac1{24}\le w_J\le\frac1{12}.
\end{equation}
Direct summation, or cancellation of the sign characters, gives
\begin{equation}\label{eq:wmoments}
 \sum_{J\ni i}w_J=1,\qquad
 \sum_{J\ne\varnothing}w_J=\frac{31}{16}.
\end{equation}
For $|B|\ge2$, define
\begin{equation}\label{eq:dprofile}
 d_B=1-\sum_{J\supseteq B}w_J.
\end{equation}
The pair values are
\begin{equation}\label{eq:pairprofile}
 d_{ij}=\frac12\quad(1\le i<j\le4),\qquad
 (d_{15},d_{25},d_{35},d_{45})
 =\frac1{24}(11,13,13,11).
\end{equation}

\begin{lemma}[Leading cancellation]\label{lem:leading}
Let $u=(2,1,-1,-2,0)$.
With $z_i=tu_i+O(1)$ and $D_B=t^2d_B+O(1)$,
\begin{equation}\label{eq:leading}
 t^{-5}F(tX)\longrightarrow X^5
\end{equation}
coefficientwise as $t\longrightarrow\infty$.
\end{lemma}
\begin{proof}
A singleton factor in \eqref{eq:FR} has order $t$ and a nonsingleton
factor has order at most $t^2$. A partition contributes to order $t^5$
only when every block has size one or two. Hence the limit is
\[
 M(X)=\sum_{\mathcal M\text{ a matching of }K_5}
       \prod_{\{i,j\}\in\mathcal M}d_{ij}
       \prod_{k\text{ unmatched}}(X+u_k).
\]
Reflection interchanging $1$ with $4$ and $2$ with $3$, and fixing $5$,
preserves the weights and negates $u$. Thus $M$ is odd.
It is monic, and its $X^3$ coefficient is
$e_2(u)+\sum d_{ij}=-5+5=0$.

If all ten weights were $1/2$, its $X$ coefficient would be
$4-15/2+15/4=1/4$.
Only the four weights incident to $5$ are changed in
\eqref{eq:pairprofile}, by $-1/24$ on rows $1,4$ and $+1/24$ on rows
$2,3$. The $X$ coefficient of the derivative for the edge $i5$ is
$u_i^2-7/2$. Since a matching contains at most one such edge, the change
is exactly
\[
 2\left(-\frac1{24}\right)\left(4-\frac72\right)
 +2\left(\frac1{24}\right)\left(1-\frac72\right)=-\frac14.
\]
The remaining $X$ coefficient is therefore zero, and $M(X)=X^5$.
\end{proof}

\section{An integer baseline and three coefficient matrices}\label{sec:matrices}

Fix the five input coefficients. For the moment let $t$ be any positive
multiple of $96$, and set
\begin{equation}\label{eq:baseline}
 L=t^2,\qquad
 z=(2t,\ t,\ -t,\ 1-2t,\ a+9).
\end{equation}
Initially put $D_B=t^2d_B$ for every $|B|\ge2$, and add $1$ to
$D_{23}$ and $D_{34}$. All these parameters are integers. In fact
$t^2/192$ is a multiple of $24$, so every quadratic profile term is a
multiple of $24$.
The trace identity \eqref{eq:trace} gives $[x^4]F=a$ exactly.

We change only parameters whose subsets contain $5$. Throughout the
construction, all outer data remain those just specified.
All coefficient vectors in the mathematical construction below are in
\emph{descending} degree order.

\subsection{Four star-pair variables}

Use $D_{15},D_{25},D_{35},D_{45}$. Their derivative coefficient matrix,
with rows $(x^3,x^2,x,1)$, is
\begin{equation}\label{eq:A}
 A(t)=
 \begin{pmatrix}
 1&1&1&1\\
 -2t-2&-t-2&t-2&2t-3\\
 \frac12t^2+2t+3&-\frac52t^2+2t+2&-\frac52t^2+2t+1&
                    \frac12t^2-2t+3\\
 p_1&p_2&p_3&p_4
 \end{pmatrix},
\end{equation}
where
\begin{align*}
 p_1&=t^3-2t^2-t+1,&
 p_2&=\frac72t^3-3t^2+2t,\\
 p_3&=-\frac72t^3+t^2,&
 p_4&=-t^3-\frac32t^2+2t.
\end{align*}
For example, for a three-element outer subset $R$ its palette polynomial is
\begin{align*}
 F_R(x)={}&x^3+\left(\sum_{i\in R}z_i-3\right)x^2\\
 &+\left(e_2(z_R)-\sum_{i\in R}z_i+\sum_{\{i,j\}\subset R}D_{ij}+2\right)x\\
 &+e_3(z_R)+\sum_{\{i,j\}\subset R}D_{ij}z_{R\setminus\{i,j\}}-2D_R.
\end{align*}
Here $z_{\{k\}}$ in the last term means $z_k$.
Substitution gives \eqref{eq:A}. Its determinant factors as
\begin{equation}\label{eq:detA}
 \det A=(3t-1)(2t^2-2t+1)(12t^3-3t^2+5t-1)>0
 \quad(t\ge1).
\end{equation}
Modulo $24$, the same matrix is
\begin{equation}\label{eq:Amod}
 A\equiv\overline A:=
 \begin{pmatrix}
 1&1&1&1\\
 -2&-2&-2&-3\\
 3&2&1&3\\
 1&0&0&0
 \end{pmatrix}\pmod{24},
 \qquad \det\overline A=1.
\end{equation}
Thus its coefficient map is invertible over $\Z/24\Z$.

\subsection{Three triple variables and two quadruple variables}

Use the triples $345,245,125$, whose complementary outer pairs are
$12,13,34$. Their derivatives are $-2F_{12},-2F_{13},-2F_{34}$.
Since $F_{ij}=(x+z_i)(x+z_j)-(x-D_{ij})$, their matrix in degrees
$(x^2,x,1)$ is
\begin{equation}\label{eq:B}
 \begin{aligned}
 B(t)&=
 \begin{pmatrix}
 -2&-2&-2\\
 2-6t&2-2t&6t\\
 -5t^2&3t^2&-5t^2+2t-2
 \end{pmatrix},\\
 \det B&=16t(12t^2-3t+1)>0.
 \end{aligned}
\end{equation}
Finally, use $D_{2345},D_{1345}$, complementary to rows $1,2$.
Their derivatives are $6(x+2t),6(x+t)$, with matrix
\begin{equation}\label{eq:C}
 C(t)=\begin{pmatrix}6&6\\12t&6t\end{pmatrix},
 \qquad \det C=-36t\ne0.
\end{equation}
The remaining full-set derivative is
$\partial F/\partial D_{12345}=-24$.

\section{Exact integral correction}\label{sec:lift}

Every change in this section is added to the indicated deficit $D_B$.
The row-affine identity, Lemma~\ref{lem:affine}, makes the effect of all
changes exactly the sum of their derivative polynomials.

\begin{proposition}[The integral lift]\label{prop:lift}
For every input $(a,b,c,d,e)\in\Z^5$ and every positive multiple $t$ of $96$,
the following finite procedure produces integer parameters with
\begin{equation}\label{eq:Fisf}
 F(x)=f(x).
\end{equation}
For each fixed input, all changes to the initial deficits are $O_f(t)$.
\end{proposition}
\begin{proof}
The monic and degree-four coefficients already agree. Let
$r=(r_3,r_2,r_1,r_0)^T$ be the remaining coefficients of $f-F$.

\paragraph{Star pairs.}
Solve $A v=r$ over $\Q$, and let $v^\flat$ be its componentwise floor.
After applying $v^\flat$ to the four star variables, the residual is
$r'=r-A v^\flat$.
Choose the unique residue vector
$\rho\in\{0,\ldots,23\}^4$ satisfying
$\overline A\rho\equiv r'\pmod{24}$, and apply it.
Now every residual coefficient is divisible by $24$.
Let $s_3$ be the current coefficient of $x^3$ in $f-F$.
Apply the additional change $s_3$ to $D_{15}$.
Its derivative has leading coefficient $1$, so the cubic coefficient
now agrees exactly. This last change is a multiple of $24$ and
preserves divisibility of all residual coefficients. Write the residual as
\[
 R_2x^2+R_1x+R_0,\qquad R_2,R_1,R_0\in24\Z.
\]

\paragraph{Triples.}
Solve
$Bv=(R_2,R_1,R_0)^T$ over $\Q$. Put
\[
 \tau_1=12\left\lfloor\frac{v_1}{12}\right\rfloor,\qquad
 \tau_2=12\left\lfloor\frac{v_2}{12}\right\rfloor,\qquad
 \tau_3=-\frac{R_2}{2}-\tau_1-\tau_2.
\]
All three integers lie in $12\Z$.
Apply these changes to $D_{345},D_{245},D_{125}$.
Their sum is the exact sum of the real solution, so the $x^2$
coefficient is removed. Since the derivative coefficients are multiples
of $2$, these changes preserve divisibility of the remaining coefficients
by $24$. Write the new residual as $S_1x+S_0$, with $S_1,S_0\in24\Z$.

\paragraph{Quadruples.}
Solve $Cv=(S_1,S_0)^T$ over $\Q$, and set
\[
 \kappa_1=4\left\lfloor\frac{v_1}{4}\right\rfloor,\qquad
 \kappa_2=\frac{S_1}{6}-\kappa_1.
\]
Both changes belong to $4\Z$. Apply them to $D_{2345},D_{1345}$.
The linear coefficient is removed, and the remaining constant
$V=f-F$ is divisible by $24$.

\paragraph{The full set.}
Apply the integer change $-V/24$ to $D_{12345}$.
Its derivative is $-24$, so the final residual is zero.

All inverse matrices exist by \eqref{eq:detA}--\eqref{eq:C}.
The bounded-size statements follow either from the estimates in
Section~\ref{sec:effective}, or from the following scaling argument.
The initial residual has orders
$(O_f(t),O_f(t^2),O_f(t^3),O_f(t^4))$.
The matrix $A$ is
$\diag(1,t,t^2,t^3)(M+O(t^{-1}))$, where
\[
 M=\begin{pmatrix}
 1&1&1&1\\-2&-1&1&2\\
 1/2&-5/2&-5/2&1/2\\
 1&7/2&-7/2&-1
 \end{pmatrix},\qquad \det M=72.
\]
Hence its exact rational correction is $O_f(t)$.
Rounding and the residue repair have bounded error, leaving residual
orders $(O(t),O(t^2),O(t^3))$ after the cubic coefficient is fixed.
The scaled triple matrix has a nonsingular limit, so its correction
is $O(t)$. The two floored coordinates have bounded error, and fixing
the sum gives bounded error also in the third coordinate. The new residual
has orders $(O(t),O(t^2))$.
The same argument for $C$ gives $O(t)$ corrections, and the final constant
residual is $O(t)$. Thus the full-set correction is $O(t)$ as well.
\end{proof}

\begin{proof}[Existence in Theorem~\ref{thm:main}]
The initial intersection moments differ from those of the atoms
$w_Jt^2$ by $O_f(t)$, and Proposition~\ref{prop:lift} preserves that
property. Möbius inversion therefore gives
\[
 m_J=w_Jt^2+O_f(t),\qquad
 N=\frac{31}{16}t^2+O_f(t),\qquad
 n:=N-L=\frac{15}{16}t^2+O_f(t).
\]
Since all $w_J\ge1/24$, for every fixed input all sufficiently large
multiples of $96$ give positive integer atoms and $n>0$.
These atoms construct the actual graph. Equations \eqref{eq:count}
and \eqref{eq:Fisf} prove \eqref{eq:main}.

The five positive singleton atoms supply distinct old colors to
$v_1,\ldots,v_5$, so $G$ is $N$-colorable; its $N$-clique shows
$\chi(G)=N$. A large-clique vertex of type $\{1\}$ is adjacent to $v_2$,
joining the two cliques, so $G$ is connected.
The next section proves that the fixed choice \eqref{eq:scale} is sufficient
and establishes the quantitative bounds \eqref{eq:bounds}.
\end{proof}

\section{A predetermined scale and uniform height bound}\label{sec:effective}

The existence proof above already treats every integer input.
We now give explicit constants, so the construction need not search for
a sufficiently large $t$.

\begin{proposition}\label{prop:effective}
With \eqref{eq:scale}, every atom satisfies
\begin{equation}\label{eq:atomexplicit}
 m_J\ge\frac{t^2}{96}>0.
\end{equation}
Furthermore,
\begin{equation}\label{eq:Nnexplicit}
 \frac{29}{32}t^2\le n\le\frac{31}{32}t^2,\qquad
 N+5\le2t^2.
\end{equation}
\end{proposition}
\begin{proof}
All estimates before the final scale choice hold for $t\ge1$ whenever
the baseline and residue steps are integral. Denote the baseline
coefficients in degrees $3,2,1,0$ by $F_3,F_2,F_1,F_0$.
Expanding \eqref{eq:FR} gives
\begin{align}
 F_3={}&-5a+2t-14,\notag\\
 F_2={}&-2at^2+2at+10a-\frac{743}{24}t^2+15t+46,\notag\\
 F_1={}&-\frac72at^2+at-5a+\frac76t^3+\frac{61}{24}t^2+12t-24,\label{eq:Fbaseline}\\
 F_0={}&\frac94at^4-at^3+\frac{11}{4}at^2+2at
        +\frac{333}{16}t^4-\frac{85}{8}t^3
        +\frac{245}{48}t^2+18t.\notag
\end{align}
For $r=(b-F_3,c-F_2,d-F_1,e-F_0)^T$, the triangle inequality,
$t\ge1$, and $\htf\ge10$ give
\begin{equation}\label{eq:rbound}
 \left\|\diag(1,t^{-1},t^{-2},t^{-3})r\right\|_\infty
 \le14\htf t.
\end{equation}
For instance, the quadratic residual divided by $t$ is at most
\[
 \left(14|a|+|c|+\frac{743}{24}+61\right)t
 \le14\htf t.
\]
The other three coordinates have bounds $7\htf t,10\htf t,8\htf t$,
respectively.

\paragraph{The first inverse and its rounding error.}
The factors in \eqref{eq:detA} give $\det A\ge18t^6$ for $t\ge1$.
For $\widetilde A=\diag(1,t^{-1},t^{-2},t^{-3})A$,
the coefficient sums of the scaled adjugate entries are bounded by
\[
 K_A=\begin{pmatrix}
 380&110&40&20\\
 150&101/2&38&37\\
 91&48&36&24\\
 109&119/2&28&7
 \end{pmatrix}.
\]
More precisely, the polynomial $(\operatorname{adj}A)_{ij}$ has degree
at most $6-(j-1)$ and coefficient absolute sum $(K_A)_{ij}$ or less.
After multiplication of column $j$ by $t^{j-1}$ and division by
$\det A$, this yields
\begin{equation}\label{eq:Ainversebound}
 \|\widetilde A^{-1}\|_\infty\le\frac{550}{18}=\frac{275}{9}.
\end{equation}
These finite coefficient bounds follow by expanding the explicit matrix
\eqref{eq:A}. Therefore every exact rational star correction has magnitude
at most $428\htf t$.

After flooring and the residue adjustment, its error relative to the
rational solution is in $(-1,23]$ in each coordinate.
The last cubic repair makes the first error equal to minus the sum of
the other three. The error vector consequently has zero sum and
$\ell^1$ norm at most $138$.
The remaining residual satisfies
\begin{equation}\label{eq:rtriplebound}
 |R_2|\le690t,\qquad |R_1|\le897t^2,\qquad |R_0|\le1173t^3.
\end{equation}
For example, the coefficient absolute bounds for a column of $A$ in these
three rows are $5t,(13/2)t^2,(17/2)t^3$; multiplication by $138$ gives
the displayed bounds. This residual is independent of the size of the
original real correction.

\paragraph{Triples.}
For $\widetilde B=\diag(1,t^{-1},t^{-2})B$, we have
$\det B\ge160t^3$ and the corresponding adjugate coefficient-sum bounds
\[
 K_B=\begin{pmatrix}34&24&20\\102&8&28\\44&16&8\end{pmatrix}.
\]
The degrees in its columns are at most $3,2,1$, respectively.
Thus
\[
 \|\widetilde B^{-1}\|_\infty\le\frac{138}{160}=\frac{69}{80}.
\]
Equation~\eqref{eq:rtriplebound} bounds each rational triple correction
by $1012t$.
Rounding two coordinates downward to multiples of $12$ and preserving
their sum gives a rounding error of $\ell^1$ norm less than $48$.
The total magnitude of the three actual changes is at most $3036t+48$.
The new residual obeys
\begin{equation}\label{eq:rquadbound}
 |S_1|\le384t,\qquad |S_0|\le432t^2,
\end{equation}
by applying that rounding-error bound to the last two rows of $B$.

\paragraph{Quadruples and the full set.}
The inverse of $\diag(1,t^{-1})C$ has infinity norm $1/2$.
Hence every rational quadruple correction has magnitude at most $216t$.
The two integer corrections have total magnitude at most $432t+8$.
Their rounding errors are opposite and have magnitude less than $4$
in each coordinate. Since $z_1-z_2=t$, the remaining constant has
magnitude less than $24t$, and the final full-set deficit correction has
magnitude less than $t$.

\paragraph{Atoms and the shift.}
Relative to the unperturbed moments
$t^2\sum_{J\supseteq B}w_J$, the singleton changes have total magnitude
at most $6t+\htf-2$, and the two initial outer-pair offsets contribute $2$.
The preceding estimates give the following bound for the sum of the
absolute changes of \emph{all thirty-one moments}:
\begin{equation}\label{eq:totalerror}
 1712\htf t+3475t+\htf+194
 <2080\htf t<3000\htf t.
\end{equation}
Each atom in \eqref{eq:mobius}, and also $N$ in \eqref{eq:N}, is an
alternating sum of a subset of these moments. Therefore
\[
 |m_J-w_Jt^2|\le3000\htf t,\qquad
 \left|N-\frac{31}{16}t^2\right|\le3000\htf t.
\]
At $t=96000\htf$, the right-hand error is $t^2/32$, and hence
\[
 m_J\ge\frac{t^2}{24}-\frac{t^2}{32}=\frac{t^2}{96},
 \qquad
 \left|n-\frac{15}{16}t^2\right|\le\frac{t^2}{32}.
\]
This proves the atom and shift bounds.
Finally,
$N+5\le(63/32)t^2+5\le2t^2$, since $t\ge96$.
Proposition~\ref{prop:effective}, and therefore Theorem~\ref{thm:main},
are proved.
\end{proof}

\section{Arithmetic consequences and algorithmic form}\label{sec:corollaries}

\begin{corollary}[Root fields and exact splitting fields]\label{cor:fields}
Every number field of degree at most five is generated by a chromatic root.
For every nonzero polynomial $g\in\Q[x]$ of degree at most five, its splitting field
over $\Q$ is exactly the splitting field of the chromatic polynomial of
a connected finite simple graph.
\end{corollary}
\begin{proof}
A number field has an integral primitive element $\alpha$.
Corollary~\ref{cor:degree} gives a chromatic root $\alpha+n$ and
$\Q(\alpha+n)=\Q(\alpha)$.

For the second statement, first make $g$ monic, let $j=\deg g$, and choose
a positive integer $m$ clearing its coefficient denominators.
Then $m^j g(x/m)$ is monic and integral and has the same splitting field
as $g$. Multiply it by $x^{5-j}$ and apply Theorem~\ref{thm:main}.
Integer translation preserves the splitting field, and all zeros of the
extra falling-factorial factor are rational. Thus equality of splitting
fields follows from \eqref{eq:main}.
For constant polynomials the splitting field is $\Q$, which is realized
by a complete graph.
\end{proof}

This proves the degree-five instance ($j=5$) of the splitting-field
realization statement in \cite[Conjecture 7.1]{Gilmore}.
It concerns the splitting field of a polynomial of degree $j$; the
splitting field itself can have larger degree.
In particular, a nonnormal quintic field is asserted to occur as a
\emph{root field}, not to be itself a splitting field.

\begin{corollary}[Bounded arithmetic description]\label{cor:algorithm}
There is an algorithm that takes five integers and returns $n,N$ and
thirty-one atom counts using a bounded number of exact rational and
integer arithmetic operations. With scale \eqref{eq:scale}, every
integer numerator, denominator, and output in this computation has
$O(\log\htf)$ bits.
\end{corollary}
\begin{proof}
All matrices have fixed size, and their entries are fixed-degree
polynomials in $t$ and the input. The only divisions are those in three
nonsingular rational systems, the fixed residue inverse, and the
displayed floor operations. At $t=96000\htf$, each resulting numerator
and denominator has polynomially bounded magnitude in $\htf$.
The graph is specified by its two cliques and the thirty-one atom
multiplicities.
\end{proof}

The bounded operation count concerns this compressed description.
Listing all vertices and edges has its own output cost.
For smaller examples, one may instead try $t=96,192,384,\ldots$ and stop
at the first positive certificate. The existence proof guarantees
termination, but the theorem itself uses the predetermined scale.

\section{An irreducible quintic example and independent counting}
\label{sec:example}

Take $f(x)=x^5-2$, which is irreducible by Eisenstein's criterion at $2$.
The construction is already valid at
\[
 t=1536,\qquad L=2\,359\,296,\qquad
 N=4\,510\,002,\qquad n=2\,150\,706.
\]
To specify the entire graph, encode a nonempty type $J$ by the integer
$s(J)=\sum_{i\in J}2^{i-1}$. Thus type $19$ is $\{1,2,5\}$.
The thirty-one atom counts are:
\begin{center}
\begin{tabular}{cr@{\qquad}cr@{\qquad}cr}
\toprule
$s(J)$ & $m_J$ & $s(J)$ & $m_J$ & $s(J)$ & $m_J$\\
\midrule
1 & 72\,406 & 12 & 152\,623 & 23 & 122\,604\\
2 & 132\,031 & 13 & 100\,548 & 24 & 202\,058\\
3 & 165\,600 & 14 & 184\,732 & 25 & 174\,276\\
4 & 166\,742 & 15 & 147\,180 & 26 & 125\,656\\
5 & 169\,788 & 16 & 86\,319 & 27 & 171\,756\\
6 & 183\,907 & 17 & 225\,578 & 28 & 142\,288\\
7 & 172\,308 & 18 & 164\,418 & 29 & 194\,364\\
8 & 89\,784 & 19 & 129\,312 & 30 & 110\,180\\
9 & 120\,636 & 20 & 126\,636 & 31 & 147\,732\\
10 & 169\,256 & 21 & 125\,124 &  & \\
11 & 123\,156 & 22 & 111\,004 &  & \\
\bottomrule
\end{tabular}
\end{center}
Every entry is positive; the minimum is $72\,406$.
Their sum is $N$. The graph has $4\,510\,007$ vertices and satisfies
\begin{equation}\label{eq:example}
 P_G(q)=\fall q{4\,510\,002}
       \bigl((q-2\,150\,706)^5-2\bigr).
\end{equation}
In particular, $2^{1/5}+2\,150\,706$ is a chromatic root.
No minimality claim is made for the graph order or the shift.

A check independent of the partition-coefficient calculation is available.
Let $r_k$ count the $k$-edge matchings of the bipartite complement of $G$.
A color class of $G$ contains at most one vertex from each of its two
cliques, so matching color classes gives the standard expansion
\begin{equation}\label{eq:rook}
 P_G(q)=\sum_{k=0}^5r_k\fall q{N+5-k}
       =\fall qN\sum_{k=0}^5r_k\fall{q-N}{5-k}.
\end{equation}
This is the biclique matching formula \cite[Section 2.1]{Bohn}.
Direct injection counting on the atom types gives, for the example,
\begin{align*}
 r_0&=1,\\
 r_1&=11\,796\,490,\\
 r_2&=55\,662\,846\,935\,065,\\
 r_3&=131\,325\,132\,122\,534\,707\,215,\\
 r_4&=154\,917\,363\,795\,559\,508\,146\,913\,281,\\
 r_5&=73\,099\,121\,379\,762\,419\,852\,518\,386\,302\,974.
\end{align*}
Substitution in \eqref{eq:rook} gives $(y+L)^5-2$ exactly for $y=q-N$.

Three independently written implementations were compared within this work.
One separate matching routine checked six selected inputs and 128
deterministic-seed random inputs with coefficient magnitudes up to $10^{24}$.
Another implementation checked four selected inputs and 100 random inputs
in $[-100,100]^5$. The implementation below was also checked at the fixed
scale \eqref{eq:scale}, including mixed coefficients of magnitude $10^{25}$.
These are implementation checks. The proof of universality is the
construction and the inequalities above, not an inference from sampling.

\section{Scope, provenance, and priority}\label{sec:scope}

The result addresses every monic integer quintic and therefore the original
$\alpha+n$ assertion for all algebraic integers of degree at most five.
The proof has no conditional conjectural input. It does not establish the
degree-six or arbitrary-degree case.

The partition and matching identities, the biclique framework, and Boolean
inversion are standard parts of the prior approach.
The specific contribution presented here is the combination of the strict
rational profile, the invertible residue matrix, the layered integral
correction, and the uniform estimates that give a genuine graph for every
quintic input. These estimates supply the existence step, beyond a
coordinate transformation or a system of necessary equations.

The supplied VMRC corpus was inspected at the following commit:
\begin{center}
\texttt{9f8547bfd3b8f6560f20c1b63028dc87a07a2e33}.
\end{center}
Its manuscript 245 informed the proof-design principle of controlling
correction costs inside a positive realization margin. The present
chromatic proof does not assume that manuscript's independence-polynomial
results, or any earlier quartic theorem, as established lemmas.

The literature assessment consulted the full texts of
\cite{CM,Bohn,Gilmore,DM22}, the official abstract of \cite{DM19}, and
Cameron's conjecture list dated 2 November 2026 \cite{CameronWeb}.
The last list concerns the general conjecture and does not by itself
certify the separate status of degree five.
The full text of \cite{DM19} was not available during preparation, and
its intermediate constructions have not been compared line by line.
No identical unrestricted quintic solution was identified in the
consulted material. Absence of an identified prior result is not a proof
of absolute publication priority. The manuscript has not undergone
external peer review or proof-assistant formalization.

\appendix
\section{Self-contained exact-arithmetic implementation}\label{app:code}

The following code uses Python 3.10 or later and only the standard library.
All polynomial arrays in the code are in \emph{ascending} coefficient
order, while matrices in the mathematical text use descending order.
Row indices in the code are $0,\ldots,4$, with row $4$ distinguished.
Integer masks encode types as in Section~\ref{sec:example}.

The function \texttt{construct\_bounded} uses the proved scale and performs
no search. The optional \texttt{construct} tries smaller doubling scales.
The independent \texttt{rook\_numbers} function counts injections into atom
types and does not use the partition formula. It verifies the full
polynomial in \eqref{eq:rook} with exact integers.

\begin{lstlisting}
"""Exact quintic palette construction, standard-library implementation.

Rows 0,1,2,3 are fixed; row 4 is the extension row.
All polynomial coefficient arrays are ascending.
"""
from fractions import Fraction as Q
from itertools import combinations
from math import factorial, comb
import json

ALL=31
OUTER=(0,1,2,3)
ROWS=(0,1,2,3,4)
SLOPE=(2,1,-1,-2,0)

def multiply(p,q):
    out=[0]*(len(p)+len(q)-1)
    for i,a in enumerate(p):
        for j,b in enumerate(q): out[i+j]+=a*b
    return out

def partitions(rows):
    if not rows:
        yield ()
        return
    first,*tail=rows
    for p in partitions(tail):
        yield (1<<first,)+p
        for j in range(len(p)):
            yield p[:j]+(p[j]|(1<<first),)+p[j+1:]

PARTS={m:tuple(partitions([i for i in ROWS if m>>i&1])) for m in range(32)}

def poly(z,D,mask=ALL):
    out=[0]*(mask.bit_count()+1)
    for p in PARTS[mask]:
        term=[1]
        for block in p:
            k=block.bit_count()
            mu=(-1)**(k-1)*factorial(k-1)
            value=z[block.bit_length()-1] if k==1 else -D[block]
            term=multiply(term,[mu*value,mu])
        for i,v in enumerate(term): out[i]+=v
    return out

def derivative(z,D,block):
    k=block.bit_count()
    scalar=(-1)**k*factorial(k-1)
    return [scalar*v for v in poly(z,D,ALL^block)]

def solve(columns,rhs):
    n=len(rhs)
    M=[[Q(columns[j][i]) for j in range(n)]+[Q(rhs[i])] for i in range(n)]
    for j in range(n):
        pivot=next(i for i in range(j,n) if M[i][j])
        M[j],M[pivot]=M[pivot],M[j]
        v=M[j][j]
        M[j]=[x/v for x in M[j]]
        for i in range(n):
            if i==j: continue
            v=M[i][j]
            M[i]=[a-v*b for a,b in zip(M[i],M[j])]
    return [M[i][-1] for i in range(n)]

def leading_atoms():
    result={}
    for m in range(1,32):
        s=[1 if m>>i&1 else -1 for i in ROWS]
        result[m]=Q(1,16)+Q(s[4]*(s[0]+s[3]-s[1]-s[2]),192)
    return result

WEIGHT=leading_atoms()
BASE={m:sum(v for J,v in WEIGHT.items() if J&m==m) for m in range(1,32)}
A0_COLUMNS=((1,3,-2,1),(0,2,-2,1),(0,1,-2,1),(0,3,-3,1))

def initial(coeff,t):
    assert t>0 and t%96==0
    a,b,c,d,e=coeff
    z=[2*t,t,-t,-2*t+1,a+9]
    D={m:Q((1-v)*t*t) for m,v in BASE.items() if m.bit_count()>1}
    D[6]+=1                 # outer pair {2,3}
    D[12]+=1                # outer pair {3,4}
    assert all(v.denominator==1 for v in D.values())
    return z,{m:int(v) for m,v in D.items()}

def atoms_from(z,D,L):
    intersections={m:L+z[m.bit_length()-1] if m.bit_count()==1 else L-D[m]
                   for m in range(1,32)}
    return {m:sum((-1)**((J^m).bit_count())*v for J,v in intersections.items() if J&m==m)
            for m in range(1,32)}

def construct_at(coeff,t):
    a,b,c,d,e=coeff
    target=[e,d,c,b,a,1]
    z,D=initial(coeff,t)
    stage_changes={}
    def residual():
        return [f-g for f,g in zip(target,poly(z,D))]
    def change(m,delta,stage):
        assert Q(delta).denominator==1
        D[m]+=int(delta)
        stage_changes.setdefault(stage,[]).append(int(delta))
    r=residual()
    assert r[4:]==[0,0]

    stars=[16|(1<<i) for i in OUTER]
    A=[derivative(z,D,m) for m in stars]
    assert [[v%24 for v in p] for p in A]==[[v%24 for v in p] for p in A0_COLUMNS]
    solution=solve(A,r[:4])
    for m,v in zip(stars,solution): change(m,v//1,'pairs')
    r=residual()
    residue=solve(A0_COLUMNS,r[:4])
    assert all(v.denominator==1 for v in residue)
    for m,v in zip(stars,residue): change(m,int(v)%24,'pairs')
    r=residual()
    assert all(v%24==0 for v in r)
    change(stars[0],r[3],'pairs')
    r=residual()
    assert r[3:]==[0,0,0]

    # The complementary outer pairs are {1,2}, {1,3}, {3,4}.
    triples=[28,26,19]
    B=[derivative(z,D,m) for m in triples]
    solution=solve(B,r[:3])
    corrections=[12*(v//12) for v in solution[:2]]
    corrections.append(-r[2]//2-sum(corrections))
    for m,v in zip(triples,corrections): change(m,v,'triples')
    r=residual()
    assert r[2:]==[0,0,0,0] and all(v%24==0 for v in r)

    quads=[30,29]
    C=[derivative(z,D,m) for m in quads]
    solution=solve(C,r[:2])
    corrections=[4*(solution[0]//4)]
    corrections.append(r[1]//6-corrections[0])
    for m,v in zip(quads,corrections): change(m,v,'quadruples')
    r=residual()
    assert r[1:]==[0,0,0,0,0] and r[0]%24==0
    change(31,-r[0]//24,'quintuple')
    assert poly(z,D)==target

    atoms=atoms_from(z,D,t*t)
    N=sum(atoms.values())
    result=dict(coefficients=[1,a,b,c,d,e],t=t,L=t*t,z=z,D=D,atoms=atoms,
                N=N,n=N-t*t,minimum_atom=min(atoms.values()),stage_changes=stage_changes)
    result['valid']=result['minimum_atom']>0 and result['n']>0
    return result

def rook_numbers(atoms):
    dp=[0]*32
    dp[0]=1
    for mask,count in atoms.items():
        mask=int(mask)
        assert count>=0
        falling=[1]
        for k in range(1,6): falling.append(falling[-1]*(count-k+1))
        nxt=[0]*32
        for used,ways in enumerate(dp):
            if not ways: continue
            available=mask & (31^used)
            sub=available
            while True:
                nxt[used|sub]+=ways*falling[sub.bit_count()]
                if sub==0: break
                sub=(sub-1)&available
        dp=nxt
    return [sum(v for m,v in enumerate(dp) if m.bit_count()==k) for k in range(6)]

def verify(cert):
    assert cert['valid']
    rook=rook_numbers(cert['atoms'])
    actual=[0]*6
    for k,r in enumerate(rook):
        p=[1]
        for j in range(5-k): p=multiply(p,[-j,1])
        for j,v in enumerate(p): actual[j]+=r*v
    target=[0]*6
    descending=cert['coefficients']
    L=cert['L']
    for degree,coefficient in enumerate(reversed(descending)):
        for j in range(degree+1):
            target[j]+=coefficient*comb(degree,j)*L**(degree-j)
    assert actual==target,(actual,target)
    return rook

def construct(coeff):
    t=96
    while True:
        result=construct_at(coeff,t)
        if result['valid']:
            result['rook_numbers']=verify(result)
            return result
        t*=2

def construct_bounded(coeff):
    height=10+sum(abs(v) for v in coeff)
    result=construct_at(coeff,96000*height)
    assert result['minimum_atom']>=result['t']**2//96
    assert 0<result['n']<result['t']**2
    assert result['N']+5<=2*result['t']**2
    result['rook_numbers']=verify(result)
    return result

if __name__ == '__main__':
    example=construct((0,0,0,0,-2))
    print(json.dumps(example,indent=2,sort_keys=True))
    bounded=construct_bounded((0,0,0,0,-2))
    print('Uniform scale verified:', bounded['N']+5)
\end{lstlisting}

\begin{thebibliography}{9}
\bibitem{CM}
P.~J. Cameron and K. Morgan,
\emph{Algebraic properties of chromatic roots},
Electronic Journal of Combinatorics \textbf{24}(1) (2017), P1.21.
\href{https://doi.org/10.37236/6578}{doi:10.37236/6578}.
\href{https://emis.de/ft/29131}{Full text}.

\bibitem{Bohn}
A. Bohn, \emph{Chromatic roots as algebraic integers},
Discrete Mathematics and Theoretical Computer Science Proceedings
\textbf{AR} (2012), 543--554.
\href{https://doi.org/10.46298/dmtcs.3061}{doi:10.46298/dmtcs.3061}.
\href{https://www.math.nagoya-u.ac.jp/fpsac12/download/contributed/dmAR0148.pdf}{Full text}.

\bibitem{Gilmore}
H.~J. Gilmore,
\emph{Algebraic Properties of Chromatic Polynomials and Their Roots},
MSc thesis, University of Waikato, 2015.
\url{https://hdl.handle.net/10289/9367}.

\bibitem{DM19}
D. Delbourgo and K. Morgan,
\emph{An algorithm which outputs a graph with a specified chromatic factor},
Discrete Applied Mathematics \textbf{257} (2019), 128--150.
\href{https://doi.org/10.1016/j.dam.2018.10.033}{doi:10.1016/j.dam.2018.10.033}.
Official abstract consulted through the author's
\href{https://profiles.auckland.ac.nz/daniel-delbourgo/publications}{University of Auckland publications record}.

\bibitem{DM22}
D. Delbourgo and K. Morgan, \emph{On the fifth chromatic coefficient},
Australasian Journal of Combinatorics \textbf{84}(1) (2022), 56--85.
\url{https://ajc.maths.uq.edu.au/pdf/84/ajc_v84_p056.pdf}.

\bibitem{VMRC}
Verifiable Mathematical Reasoning Corpus, manuscript 245,
\emph{Roots and coefficient peaks of graph polynomials},
research manuscript dated 21 November 2026, Sections 2--3.
Snapshot \texttt{9f8547bfd3b8f6560f20c1b63028dc87a07a2e33},
accessed 28 November 2026.
\url{https://github.com/QEEDDD/Verifiable-Mathematical-Reasoning-Corpus}.

\bibitem{CameronWeb}
P.~J. Cameron, \emph{Peter Cameron's conjectures},
page dated 2 November 2026, accessed 28 November 2026.
\url{https://cameroncounts.github.io/web/problems/cameronconjs.html}.
\end{thebibliography}
\end{document}


